\ifdefined\gkver\else\def\gkver{student}\fi
\documentclass[\gkver]{gaokaozhenti}
\begin{document}

\chapter{2006年上海}

\begin{enumerate}
\item
%% number: 1
%% paperName: 2006年上海高考第 1 题 4 分
%% typeId: 3
%% type: 填空题
%% chapter: 电场
%% point: 带电粒子的偏转
%% method: 无
%% score: 4
%% degree: 850
%% duplicateId: 0
%% body:
\textbf{A．}如图所示，一束 $\beta$ 粒子自下而上进入一水平方向的匀强电场后发生偏转，则电场方向向\tkanswer[1.2cm]{左}，进入电场后，$\beta$ 粒子的动能\tkanswer[1.6cm]{增加}（填“增加”、“减少”或“不变”）。

\onepicture[w=4.2cm, align=right]{\includesvg{figs/01a}}

\textbf{B．}如图所示，一束 $\beta$ 粒子自下而上进入一垂直纸面的匀强磁场后发生偏转，则磁场方向向\tkanswer[1.2cm]{内}，进入磁场后，$\beta$ 粒子的动能\tkanswer[1.6cm]{不变}（填“增加”、“减少”或“不变”）。

\onepicture[w=4.2cm, align=right]{\includesvg{figs/01b}}

%% answer:
\jdanswer{
\textbf{A．}左，增加。

\textbf{B．}内，不变。
}

%% memo:
\memoanswer{
\textbf{A．}由图可知，$\beta$ 粒子（带负电）所受电场力方向水平向右，故电场方向向左。由于电场力做正功，由动能定理可知粒子在电场中动能增加。

\textbf{B．}$\beta$ 粒子（带负电）沿纸面向左偏转，由左手定则判断出磁场方向垂直纸面向（内）里，又洛伦兹力不做功，故粒子的动能不变。
}

\item
%% number: 2
%% paperName: 2006年上海高考第 2 题 4 分
%% typeId: 3
%% type: 填空题
%% chapter: 交流电
%% point: 变压器
%% method: 无
%% score: 4
%% degree: 850
%% duplicateId: 0
%% body:
\textbf{A．}如图所示，同一平面内有两根互相平行的长直导线 1 和 2，通有大小相等、方向相反的恒定电流，a、b 两点与两导线共面，a 点在两导线的中间，与两导线的距离均为 $r$，b 点在导线 2 右侧，与导线 2 的距离也为 $r$。现测得 a 点磁感应强度的大小为 $B$，则去掉导线 1 后，b 点的磁感应强度大小为\tkanswer[1.4cm]{$\frac{B}{2}$}，方向\tkanswer[2.4cm]{（垂直纸面）向外}。

\onepicture[w=4.2cm, align=right]{\includesvg{figs/02a}}

\textbf{B．}如图所示，一理想变压器原、副线圈匝数分别为 $n_{1}$ 和 $n_{2}$，当负载电阻 $R$ 中流过的电流为 $I$ 时，原线圈中流过的电流为\tkanswer[2.0cm]{$\frac{n_{2}}{n_{1}}I$}；现减小负载电阻 $R$ 的阻值，则变压器的输入功率将\tkanswer[1.0cm]{增大}（填“增大”、“减小”或“不变”）。

\onepicture[w=5.0cm, align=right]{\includesvg{figs/02b}}

%% answer:
\jdanswer{
\textbf{A．}$\frac{B}{2}$，（垂直纸面）向外。

\textbf{B．}$\frac{n_{2}}{n_{1}}I$，增大。
}

%% memo:
\memoanswer{
\textbf{A．}由安培定则可知，1、2 两根导线在 a 点的磁感应强度大小相等、方向相同，都为 $\frac{B}{2}$，而导线 2 在 a、b 两处的磁感应强度等大反向，故去掉导线 1 后，b 点的磁感应强度大小为 $\frac{B}{2}$，方向垂直两导线所在平面向外。

\textbf{B．}由变压器的电流比可知，$\frac{I_{1}}{I}=\frac{n_{2}}{n_{1}}$，故 $I_{1}=\frac{n_{2}}{n_{1}}I$。由题意知原、副线圈的电压不变，由 $P_{2}=\frac{U_{2}^{2}}{R}$ 和 $P_{1}=P_{2}$ 可知，$R$ 减小时，则 $P_{2}$、$P_{1}$ 均增大。
}

\item
%% number: 3
%% paperName: 2006年上海高考第 3 题 4 分
%% typeId: 3
%% type: 填空题
%% chapter: 电路
%% point: 逻辑电路
%% method: 无
%% score: 4
%% degree: 850
%% duplicateId: 0
%% body:
\textbf{A．}利用光电管产生光电流的电路如图所示。电源的正极应接在\tkanswer[0.8cm]{a}端（填“a”或“b”）；若电流表读数为 $8\UuA$，则每秒从光电管阴极发射的光电子至少是\tkanswer[2.2cm]{$5\times10^{13}$}个（已知电子电量为 $1.6\times10^{-19}\UC$）。

\onepicture[w=4.5cm, align=right]{\includesvg{figs/03a}}

\textbf{B．}右图为包含某逻辑电路的一个简单电路图，L 为小灯泡。光照射电阻 $R'$ 时，其阻值将变得远小于 $R$。该逻辑电路是\tkanswer[0.8cm]{非}门电路（填“与”、“或”或“非”）。当电阻 $R'$ 受到光照时，小灯泡 L 将\tkanswer[1.6cm]{发光}（填“发光”或“不发光”）。

\onepicture[w=4.5cm, align=right]{\includesvg{figs/03b}}

%% answer:
\jdanswer{
\textbf{A．}a，$5\times10^{13}$。

\textbf{B．}非，发光。
}

%% memo:
\memoanswer{
\textbf{A．}要使电子加速，应在 A 极接高电势，故 a 端为电源正极；由 $I=\frac{ne}{t}$，得 $n=\frac{It}{e}=5\times10^{13}$。

\textbf{B．}当电阻 $R'$ 受到光照时，非门输入低电压，输出高电压，故灯泡将发光。
}

\item
%% number: 4
%% paperName: 2006年上海高考第 4 题 4 分
%% typeId: 3
%% type: 填空题
%% chapter: 直线运动
%% point: 匀变速直线运动
%% method: 图象法
%% score: 4
%% degree: 650
%% duplicateId: 0
%% body:
伽利略通过研究自由落体和物块沿光滑斜面的运动，首次发现了匀加速运动规律。伽利略假设物块沿斜面运动与物块自由下落遵从同样的法则，他在斜面上用刻度表示物块滑下的路程，并测出物块通过相应路程的时间，然后用图线表示整个运动过程，如图所示。图中 OA 表示测得的时间，矩形 OAED 的面积表示该时间内物块经过的路程，则图中 OD 的长度表示\tkanswer[1.8cm]{平均速度}。P 为 DE 的中点，连接 OP 且延长交 AE 的延长线于 B 点，则 AB 的长度表示\tkanswer[1.8cm]{末速度}。

\onepicture[w=4.5cm, align=right]{\includesvg{figs/04}}

%% answer:
\jdanswer{平均速度，末速度}

%% memo:
\memoanswer{
匀变速直线运动的 $v-t$ 图象所围面积表示位移，中位线表示平均速度，最高点的纵坐标表示末速度。
}

\item
%% number: 5
%% paperName: 2006年上海高考第 5 题 4 分
%% typeId: 3
%% type: 填空题
%% chapter: 功和能
%% point: 动能定理
%% method: 极值
%% score: 4
%% degree: 450
%% duplicateId: 0
%% body:
半径分别为 $r$ 和 $2r$ 的两个质量不计的圆盘，共轴固定连结在一起，可以绕水平轴 O 无摩擦转动，大圆盘的边缘上固定有一个质量为 $m$ 的质点，小圆盘上绕有细绳。开始时圆盘静止，质点处在水平轴 O 的正下方位置。现以水平恒力 $F$ 拉细绳，使两圆盘转动，若恒力 $F=mg$，两圆盘转过的角度 $\theta=$\tkanswer[1.4cm]{$\frac{\pi}{6}$}时，质点 $m$ 的速度最大。若圆盘转过的最大角度 $\theta=\frac{\pi}{3}$，则此时恒力 $F=$\tkanswer[1.6cm]{$\frac{3mg}{\pi}$}。

\onepicture[w=4.5cm, align=right]{\includesvg{figs/05}}

%% answer:
\jdanswer{$\frac{\pi}{6}$，$\frac{3mg}{\pi}$}

%% memo:
\memoanswer{
速度最大的位置就是力矩平衡的位置，则有 $Fr=mg\cdot2r\sin\theta$，解得 $\theta=\frac{\pi}{6}$；若圆盘转过的最大角度 $\theta=\frac{\pi}{3}$，则质点 $m$ 速度最大时 $\theta=\frac{\pi}{6}$，故可求得 $F=\frac{3mg}{\pi}$。
}

\item
%% number: 6
%% paperName: 2006年上海高考第 6 题 5 分
%% typeId: 2
%% type: 多选题
%% chapter: 光学
%% point: 光的电磁说
%% method: 无
%% score: 5
%% degree: 850
%% duplicateId: 0
%% body:
人类对光的本性的认识经历了曲折的过程。下列关于光的本性的陈述符合科学规律或历史事实的是\xzanswer{BCD}

\fourchoices[answer=BCD]
{牛顿的“微粒说”与爱因斯坦的“光子说”本质上是一样的}
{光的双缝干涉实验显示了光具有波动性}
{麦克斯韦预言了光是一种电磁波}
{光具有波粒二象性}

%% answer: BCD
%% memo:
\memoanswer{
爱因斯坦的“光子说”和牛顿的“微粒说”除了都是粒子之外，并无任何相同之处，光子不是牛顿所描绘的那种遵循经典力学运动定律的微粒，故 BCD 正确。
}

\item
%% number: 7
%% paperName: 2006年上海高考第 7 题 5 分
%% typeId: 1
%% type: 单选题
%% chapter: 原子和原子核
%% point: 原子核式结构
%% method: 无
%% score: 5
%% degree: 650
%% duplicateId: 0
%% body:
卢瑟福通过对 $\alpha$ 粒子散射实验结果的分析，提出\xzanswer{A}

\fourchoices[answer=A]
{原子的核式结构模型}
{原子核内有中子存在}
{电子是原子的组成部分}
{原子核是由质子和中子组成的}

%% answer: A
%% memo:
\memoanswer{
本题原卷及材料中未提供详解，待补充。
}

\item
%% number: 8
%% paperName: 2006年上海高考第 8 题 5 分
%% typeId: 2
%% type: 多选题
%% chapter: 电场
%% point: 电场中的图象
%% method: 无
%% score: 5
%% degree: 650
%% duplicateId: 0
%% body:
A、B 是一条电场线上的两点，若在 A 点释放一初速为零的电子，电子仅受电场力作用，并沿电场线从 A 运动到 B，其速度随时间变化的规律如图所示。设 A、B 两点的电场强度分别为 $E_{A}$、$E_{B}$，电势分别为 $\varphi_{A}$、$\varphi_{B}$，则\xzanswer{AD}

\twopicture[showlabel=false, wA=6.5cm, wB=3.4cm]{figs/08a.png}{figs/08b.png}

\fourchoices[answer=AD]
{$E_{A}=E_{B}$}
{$E_{A}<E_{B}$}
{$\varphi_{A}=\varphi_{B}$}
{$\varphi_{A}<\varphi_{B}$}

%% answer: AD
%% memo:
\memoanswer{
由图象可知，电子做匀加速直线运动，故该电场为匀强电场，即 $E_{A}=E_{B}$，电子动能增加，电势能减少，电势升高，即 $\varphi_{A}<\varphi_{B}$，故 AD 正确。
}

\item
%% number: 9
%% paperName: 2006年上海高考第 9 题 5 分
%% typeId: 1
%% type: 单选题
%% chapter: 气体性质
%% point: 气体的状态参量
%% method: 无
%% score: 5
%% degree: 650
%% duplicateId: 0
%% body:
如图所示，竖直放置的弯曲管 A 端开口，B 端封闭，密度为 $\rho$ 的液体将两段空气封闭在管内，管内液面高度差分别为 $h_{1}$、$h_{2}$ 和 $h_{3}$，则 B 端气体的压强为（已知大气压强为 $p_{0}$）\xzanswer{B}

\onepicture[w=5.5cm]{\includesvg{figs/09}}

\fourchoices[answer=B]
{$p_{0}-\rho g(h_{1}+h_{2}-h_{3})$}
{$p_{0}-\rho g(h_{1}+h_{3})$}
{$p_{0}-\rho g(h_{1}+h_{3}-h_{2})$}
{$p_{0}-\rho g(h_{1}+h_{2})$}

%% answer: B
%% memo:
\memoanswer{
由图中液面的高度关系知，$p_{0}=p_{2}+\rho gh_{3}$ 和 $p_{2}=p_{1}+\rho gh_{1}$，联立解得 $p_{1}=p_{0}-\rho g(h_{1}+h_{3})$，故 B 正确。
}

\item
%% number: 10
%% paperName: 2006年上海高考第 10 题 5 分
%% typeId: 2
%% type: 多选题
%% chapter: 机械波
%% point: 波的多解性
%% method: 无
%% score: 5
%% degree: 450
%% duplicateId: 0
%% body:
在均匀介质中选取平衡位置在同一直线上的 9 个质点，相邻两质点的距离均为 $L$，如图\ref{2006上海10a}所示。一列横波沿该直线向右传播，$t=0$ 时到达质点 1，质点 1 开始向下运动，经过时间 $\Delta t$ 第一次出现如图\ref{2006上海10b}所示的波形。则该波的\xzanswer{BC}

\twopicture[wA=6.6cm, wB=6.6cm, labelA=2006上海10a, labelB=2006上海10b]{figs/10a.png}{figs/10b.png}

\fourchoices[answer=BC]
{周期为 $\Delta t$，波长为 $8L$}
{周期为 $\frac{2}{3}\Delta t$，波长为 $8L$}
{周期为 $\frac{2}{3}\Delta t$，波速为 $\frac{12L}{\Delta t}$}
{周期为 $\Delta t$，波速为 $\frac{8L}{\Delta t}$}

%% answer: BC
%% memo:
\memoanswer{
由题中波形图可知 $\lambda=8L$，质点 1 该时刻正向上运动，而 $t=0$ 时质点 1 开始向下运动，故传播时间 $\Delta t=\left(n+\frac{1}{2}\right)T$，由题意知此时第一次出现题中所示的波形，所以 $n=1$，传播距离 $x=1.5\lambda=12L$，周期 $T=\frac{2}{3}\Delta t$，波速 $v=\frac{\lambda}{T}=\frac{x}{t}=\frac{12L}{\Delta t}$，故 BC 正确。
}

\item
%% number: 11
%% paperName: 2006年上海高考第 11 题 5 分
%% typeId: 2
%% type: 多选题
%% chapter: 电路
%% point: 动态分析
%% method: 等效替代
%% score: 5
%% degree: 450
%% duplicateId: 0
%% body:
在如图所示电路中，闭合电键 S，当滑动变阻器的滑动触头 P 向下滑动时，四个理想电表的示数都发生变化，电表的示数分别用 $I$、$U_{1}$、$U_{2}$ 和 $U_{3}$ 表示，电表示数变化量的大小分别用 $\Delta I$、$\Delta U_{1}$、$\Delta U_{2}$ 和 $\Delta U_{3}$ 表示。下列比值正确的是\xzanswer{ACD}

\onepicture[w=5.5cm]{\includesvg{figs/11}}

\fourchoices[answer=ACD]
{$\frac{U_{1}}{I}$ 不变，$\frac{\Delta U_{1}}{\Delta I}$ 不变}
{$\frac{U_{2}}{I}$ 变大，$\frac{\Delta U_{2}}{\Delta I}$ 变大}
{$\frac{U_{2}}{I}$ 变大，$\frac{\Delta U_{2}}{\Delta I}$ 不变}
{$\frac{U_{3}}{I}$ 变大，$\frac{\Delta U_{3}}{\Delta I}$ 不变}

%% answer: ACD
%% memo:
\memoanswer{
由欧姆定律知 $\frac{\Delta U_{1}}{\Delta I}=\frac{U_{1}}{I}=R_{1}$，由于 $R_{1}$ 不变，故 $\frac{U_{1}}{I}$ 不变，$\frac{\Delta U_{1}}{\Delta I}$ 不变；同理，$\frac{U_{2}}{I}=R_{2}$，由于 $R_{2}$ 变大，所以 $\frac{U_{2}}{I}$ 变大，但是 $\frac{\Delta U_{2}}{\Delta I}=\frac{\Delta I(R_{1}+r)}{\Delta I}=R_{1}+r$，所以 $\frac{\Delta U_{2}}{\Delta I}$ 不变；而 $\frac{U_{3}}{I}=R_{2}+R_{1}$，所以 $\frac{U_{3}}{I}$ 变大；由于 $\frac{\Delta U_{3}}{\Delta I}=\frac{\Delta Ir}{\Delta I}=r$，所以 $\frac{\Delta U_{3}}{\Delta I}$ 不变，故 ACD 正确。
}

\item
%% number: 12
%% paperName: 2006年上海高考第 12 题 5 分
%% typeId: 2
%% type: 多选题
%% chapter: 电磁感应
%% point: 电磁感应能量分析
%% method: 无
%% score: 5
%% degree: 450
%% duplicateId: 0
%% body:
如图所示，平行金属导轨与水平面成 $\theta$ 角，导轨与固定电阻 $R_{1}$ 和 $R_{2}$ 相连，匀强磁场垂直穿过导轨平面。有一导体棒 ab，质量为 $m$，导体棒的电阻与固定电阻 $R_{1}$ 和 $R_{2}$ 的阻值均相等，与导轨之间的动摩擦因数为 $\mu$，导体棒 ab 沿导轨向上滑动，当上滑的速度为 $v$ 时，受到安培力的大小为 $F$。此时\xzanswer{BCD}

\onepicture[w=5.5cm]{\includesvg{figs/12}}

\fourchoices[answer=BCD]
{电阻 $R_{1}$ 消耗的热功率为 $\frac{Fv}{3}$}
{电阻 $R_{2}$ 消耗的热功率为 $\frac{Fv}{6}$}
{整个装置因摩擦而消耗的热功率为 $\mu mgv\cos\theta$}
{整个装置消耗的机械功率为 $(F+\mu mg\cos\theta)v$}

%% answer: BCD
%% memo:
\memoanswer{
由法拉第电磁感应定律得 $E=BLv$，回路总电流 $I=\frac{E}{1.5R}$，安培力 $F=BIL$，所以电阻 $R$ 的功率 $P_{1}=(0.5I)^{2}R=\frac{Fv}{6}$，故 B 正确。由于摩擦力 $f=\mu mg\cos\theta$，故因摩擦而消耗的热功率为 $\mu mgv\cos\theta$，整个装置消耗的机械功率为 $(F+\mu mg\cos\theta)v$。
}

\item
%% number: 13
%% paperName: 2006年上海高考第 13 题 5 分
%% typeId: 1
%% type: 单选题
%% chapter: 曲线运动
%% point: 平抛运动
%% method: 无
%% score: 5
%% degree: 450
%% duplicateId: 0
%% body:
如图所示，一足够长的固定斜面与水平面的夹角为 $37^{\circ}$，质点 A 以初速度 $v_{1}$ 从斜面顶端水平抛出，质点 B 在斜面上距顶端 $L=15\Um$ 处同时以速度 $v_{2}$ 沿斜面向下匀速运动，经历时间 $t$ 物体 A 和物体 B 在斜面上相遇，则下列各组速度和时间中满足条件的是\xzanswer{C}（$\sin37^{\circ}=0.6$，$\cos37^{\circ}=0.8$，$g=10\Umsq$）

\onepicture[w=5.0cm]{\includesvg{figs/13}}

\fourchoices[answer=C]
{$v_{1}=16\Ums$，$v_{2}=15\Ums$，$t=3\Us$}
{$v_{1}=16\Ums$，$v_{2}=16\Ums$，$t=2\Us$}
{$v_{1}=20\Ums$，$v_{2}=20\Ums$，$t=3\Us$}
{$v_{1}=20\Ums$，$v_{2}=16\Ums$，$t=2\Us$}

%% answer: C
%% memo:
\memoanswer{
由平抛运动规律可知，$\tan\theta=\frac{\frac{1}{2}gt^{2}}{v_{1}t}$，将 $\theta=37^{\circ}$ 代入解得 $3v_{1}=20t$，故只有 C 选项满足条件。
}

\item
%% number: 14
%% paperName: 2006年上海高考第 14 题 5 分
%% typeId: 4
%% type: 实验题
%% chapter: 原子和原子核
%% point: 人工核反应
%% method: 无
%% score: 5
%% degree: 650
%% duplicateId: 0
%% body:
1919 年卢瑟福通过如图所示的实验装置，第一次完成了原子核的人工转变，并由此发现\tkanswer[1.2cm]{质子}。图中 A 为放射源发出的\tkanswer[1.0cm]{$\alpha$}粒子，B 为\tkanswer[0.8cm]{氮}气。完成该实验的下列核反应方程\tkanswer[1.6cm]{$\ce{^{4}_{2}He}$} $+$ \tkanswer[1.6cm]{$\ce{^{14}_{7}N}$} $\to$ $\ce{^{17}_{8}O}$ $+$ \tkanswer[1.6cm]{$\ce{^{1}_{1}H}$}。

\onepicture[w=6.5cm, align=right]{\includesvg{figs/14}}

%% answer:
\jdanswer{质子，$\alpha$，氮，$\ce{^{4}_{2}He}+\ce{^{14}_{7}N}\to\ce{^{17}_{8}O}+\ce{^{1}_{1}H}$}

%% memo:
\memoanswer{
$\alpha$ 粒子轰击氮核，放出同位素氧，并发现质子，是原子核的人工转变的典型例子。
}

\item
%% number: 15
%% paperName: 2006年上海高考第 15 题 6 分
%% typeId: 4
%% type: 实验题
%% chapter: 电磁感应
%% point: 验证电磁感应实验
%% method: 无
%% score: 6
%% degree: 450
%% duplicateId: 0
%% body:
在研究电磁感应现象实验中，

\begin{enumerate}
	\item
	为了能明显地观察到实验现象，请在如图所示的实验器材中，选择必要的器材，在图中用实线连接成相应的实物电路图；
	\item
	将原线圈插入副线圈中，闭合电键，副线圈中感应电流与原线圈中电流的绕行方向\tkanswer[1.6cm]{相反}（填“相同”或“相反”）；
	\item
	将原线圈拔出时，副线圈中的感应电流与原线圈中电流的绕行方向\tkanswer[1.6cm]{相同}（填“相同”或“相反”）。
\end{enumerate}

\onepicture[w=6.0cm, align=right]{figs/15.png}

%% answer:
\jdanswer{
\begin{enumerate}
	\item
	实物电路图连线如图所示。
	\item
	相反
	\item
	相同
\end{enumerate}
}

%% memo:
\memoanswer{
（1）实物电路图连线如图所示。

\onepicture[w=6.0cm]{figs/15_an.png}

（2）将原线圈插入副线圈中，闭合电键，副线圈中感应电流与原线圈中电流的绕行方向相反。

（3）将原线圈拔出时，副线圈中的感应电流与原线圈中电流的绕行方向相同。
}

\item
%% number: 16
%% paperName: 2006年上海高考第 16 题 5 分
%% typeId: 4
%% type: 实验题
%% chapter: 气体性质
%% point: 查理定律
%% method: 无
%% score: 5
%% degree: 650
%% duplicateId: 0
%% body:
为了测试某种安全阀在外界环境为一个大气压时，所能承受的最大内部压强，某同学自行设计制作了一个简易的测试装置。该装置是一个装有电加热器和温度传感器的可密闭容器。测试过程可分为如下操作步骤：

a．记录密闭容器内空气的初始温度 $t_{1}$；

b．当安全阀开始漏气时，记录容器内空气的温度 $t_{2}$；

c．用电加热器加热容器内的空气；

d．将待测安全阀安装在容器盖上；

e．盖紧装有安全阀的容器盖，将一定量空气密闭在容器内。

\begin{enumerate}
	\item
	将每一步骤前的字母按正确的操作顺序填写：\tkanswer[2.8cm]{deacb}；
	\item
	若测得的温度分别为 $t_{1}=27\UduC$，$t_{2}=87\UduC$，已知大气压强为 $1.0\times10^{5}\UPa$，则测试结果是：这个安全阀能承受的最大内部压强是\tkanswer[2.8cm]{$1.2\times10^{5}\UPa$}。
\end{enumerate}

%% answer:
\jdanswer{
\begin{enumerate}
	\item
	deacb
	\item
	$1.2\times10^{5}\UPa$（或 1.2 个大气压）
\end{enumerate}
}

%% memo:
\memoanswer{
（1）按照先安装、再密闭、再加热、最后读数的顺序，正确的操作顺序为 deacb。

（2）由理想气体状态方程得 $\frac{p_{1}}{p_{2}}=\frac{T_{1}}{T_{2}}$，将 $T_{1}=300\UK$，$T_{2}=360\UK$，$p_{1}=1.0\times10^{5}\UPa$ 代入，可以解得 $p_{2}=1.2\times10^{5}\UPa$。
}

\item
%% number: 17
%% paperName: 2006年上海高考第 17 题 7 分
%% typeId: 4
%% type: 实验题
%% chapter: 电路
%% point: 分压与分流
%% method: 无
%% score: 7
%% degree: 450
%% duplicateId: 0
%% body:
表格中所列数据是测量小灯泡 $U-I$ 关系的实验数据：

\begin{table}[htbp]
\centering
\begin{tblr}{|c|c|c|c|c|c|c|c|c|}
\hline
$U$（$\UV$） & 0.0 & 0.2 & 0.5 & 1.0 & 1.5 & 2.0 & 2.5 & 3.0 \\
\hline
$I$（$\UA$） & 0.000 & 0.050 & 0.100 & 0.150 & 0.180 & 0.195 & 0.205 & 0.215 \\
\hline
\end{tblr}
\end{table}

\onepicture[w=13.0cm]{\includesvg{figs/17a}}

\begin{enumerate}
	\item
	分析上表内实验数据可知，应选用的实验电路图是图\tkanswer[1.0cm]{甲}（填“甲”或“乙”）；
	\item
	在方格纸内画出小灯泡的 $U-I$ 曲线。分析曲线可知小灯泡的电阻随 $I$ 变大而\tkanswer[1.2cm]{变大}（填“变大”、“变小”或“不变”）；
	\item
	如图丙所示，用一个定值电阻 $R$ 和两个上述小灯泡组成串并联电路，连接到内阻不计、电动势为 $3\UV$ 的电源上。已知流过电阻 $R$ 的电流是流过灯泡 b 电流的两倍，则流过灯泡 b 的电流约为\tkanswer[1.6cm]{$0.07\UA$}。
\end{enumerate}

\onepicture[w=6.5cm, align=right]{\includesvg{figs/17b}}

%% answer:
\jdanswer{
\begin{enumerate}
	\item
	甲
	\item
	小灯泡的 $U-I$ 图线如图所示，变大
	\item
	$0.07\UA$
\end{enumerate}
}

%% memo:
\memoanswer{
（1）因小灯泡两端的电压是从零开始的，所以滑动变阻器应采用分压接法，符合图中数据的要求，故选甲。

（2）将表中数据逐一描点后用平滑曲线连接各点，如图所示，图线的斜率表示小灯泡的电阻，所以小灯泡的电阻随 $I$ 变大而变大。

\onepicture[w=6.5cm]{\includesvg{figs/17_an}}

（3）由电路的连接方式知，流过灯 a 的电流是流过灯 b 的 3 倍，这就相当于同一个灯泡两种不同的工作状态，这两种状态下，灯泡的电压之和为 $3\UV$，电流满足 3 倍的关系，有 $E=U_{a}+U_{b}$，$I_{a}=3I_{b}$，在灯泡的 $U-I$ 曲线上找这样的两种状态，可确定灯泡的电流值约为 $0.07\UA$。
}

\item
%% number: 18
%% paperName: 2006年上海高考第 18 题 7 分
%% typeId: 4
%% type: 实验题
%% chapter: 机械振动
%% point: 单摆
%% method: 无
%% score: 7
%% degree: 250
%% duplicateId: 0
%% body:
有一测量微小时间差的装置，由两个摆长略有微小差别的单摆同轴水平悬挂构成，两个单摆摆动平面前后相互平行。

\begin{enumerate}
	\item
	现测得两单摆完成 50 次全振动的时间分别为 $50.0\Us$ 和 $49.0\Us$，则两单摆的周期差 $\Delta T=$\tkanswer[1.6cm]{$0.02\Us$}；
	\item
	某同学利用此装置测量小于单摆周期的微小时间差，具体操作如下：把两摆球向右拉至相同的摆角处，先释放长摆摆球，接着再释放短摆摆球，测得短摆经过若干次全振动后，两摆恰好第一次同时同方向通过某位置，由此可得出释放两摆的微小时间差。若测得释放两摆的时间差 $\Delta t=0.165\Us$，则在短摆释放\tkanswer[1.6cm]{$8.085\Us$}（填时间）后，两摆恰好第一次同时向\tkanswer[1.0cm]{左}（填方向）通过\tkanswer[1.6cm]{平衡位置}（填位置）；
	\item
	为了能更准确地测量微小的时间差，你认为此装置还可做的改进是\tkanswer[3.4cm]{减小两单摆的摆长差等}。
\end{enumerate}

%% answer:
\jdanswer{
\begin{enumerate}
	\item
	$0.02\Us$
	\item
	$8.085\Us$，左，平衡位置
	\item
	减小两单摆的摆长差等
\end{enumerate}
}

%% memo:
\memoanswer{
（1）$\Delta T=T_{1}-T_{2}=\frac{50}{50}\Us-\frac{49}{50}\Us=0.02\Us$。

（2）先释放的是长摆，故有 $nT_{1}=nT_{2}+\Delta t$，解得 $n=8.25$，所以短摆释放的时间为 $t=nT_{2}=8.085\Us$，此时两摆同时向左经过平衡位置。
}

\item
%% number: 19
%% paperName: 2006年上海高考第 19 题 10 分
%% typeId: 6
%% type: 计算题
%% chapter: 气体性质
%% point: 理想气体状态方程
%% method: 无
%% score: 10
%% degree: 650
%% duplicateId: 0
%% body:
\textbf{A．}一轻活塞将一定质量的理想气体封闭在水平固定放置的气缸内，开始时气体体积为 $V_{0}$，温度为 $27\UduC$。在活塞上施加压力，将气体体积压缩到 $\frac{2}{3}V_{0}$，温度升高到 $57\UduC$。设大气压强 $p_{0}=1.0\times10^{5}\UPa$，活塞与气缸壁摩擦不计。

\begin{enumerate}
	\item
	求此时气体的压强；
	\item
	保持温度不变，缓慢减小施加在活塞上的压力使气体体积恢复到 $V_{0}$，求此时气体的压强。
\end{enumerate}

\textbf{B．}一轻活塞将一定质量的理想气体封闭在气缸内，初始时气体体积为 $3.0\times10^{-3}\Umc$。用 DIS 实验系统测得此时气体的温度和压强分别为 $300\UK$ 和 $1.0\times10^{5}\UPa$。推动活塞压缩气体，测得气体的温度和压强分别为 $320\UK$ 和 $1.6\times10^{5}\UPa$，活塞与气缸壁摩擦不计。

\begin{enumerate}
	\item
	求此时气体的体积；
	\item
	保持温度不变，缓慢改变作用在活塞上的力，使气体压强变为 $8.0\times10^{4}\UPa$，求此时气体的体积。
\end{enumerate}

%% answer:
\jdanswer{
\textbf{A．}
\begin{enumerate}
	\item
	$p_{2}=1.65\times10^{5}\UPa$
	\item
	$p_{3}=1.1\times10^{5}\UPa$
\end{enumerate}

\textbf{B．}
\begin{enumerate}
	\item
	$V_{2}=2.0\times10^{-3}\Umc$
	\item
	$V_{3}=4.0\times10^{-3}\Umc$
\end{enumerate}
}

%% memo:
\memoanswer{
\textbf{A．}（1）气体从状态 Ⅰ 到状态 Ⅱ 的变化符合理想气体状态方程

$\frac{p_{1}V_{1}}{T_{1}}=\frac{p_{2}V_{2}}{T_{2}}$ ①

由①式 $p_{2}=\frac{p_{1}V_{1}T_{2}}{V_{2}T_{1}}=\frac{1\times10^{5}\times V_{0}\times(273+57)}{(273+27)\times\frac{2}{3}V_{0}}=1.65\times10^{5}\UPa$ ②

（2）气体从状态 Ⅱ 到状态 Ⅲ 的变化为等温过程

$p_{2}V_{2}=p_{3}V_{3}$ ③

由③式 $p_{3}=\frac{p_{2}V_{2}}{V_{3}}=\frac{1.65\times10^{5}\times\frac{2}{3}V_{0}}{V_{0}}=1.1\times10^{5}\UPa$ ④

\textbf{B．}（1）气体从状态 Ⅰ 到状态 Ⅱ 的变化符合理想气体状态方程

$\frac{p_{1}V_{1}}{T_{1}}=\frac{p_{2}V_{2}}{T_{2}}$ ①

由①式 $V_{2}=\frac{p_{1}V_{1}T_{2}}{p_{2}T_{1}}=\frac{1\times10^{5}\times3\times10^{-3}\times320}{1.6\times10^{5}\times300}=2.0\times10^{-3}\Umc$ ②

（2）气体从状态 Ⅱ 到状态 Ⅲ 的变化为等温过程

$p_{2}V_{2}=p_{3}V_{3}$ ③

由③式 $V_{3}=\frac{p_{2}V_{2}}{p_{3}}=\frac{1.6\times10^{5}\times2.0\times10^{-3}}{8.0\times10^{4}}=4.0\times10^{-3}\Umc$ ④
}

\item
%% number: 20
%% paperName: 2006年上海高考第 20 题 10 分
%% typeId: 5
%% type: 辨析题
%% chapter: 直线运动
%% point: 多段匀变速直线运动
%% method: 无
%% score: 10
%% degree: 450
%% duplicateId: 0
%% body:
要求摩托车由静止开始在尽量短的时间内走完一段直道，然后驶入一段半圆形的弯道，但在弯道上行驶时车速不能太快，以免因离心作用而偏出车道，求摩托车在直道上行驶所用的最短时间。有关数据见表格。

\begin{table}[htbp]
\centering
\begin{tblr}{|c|c|}
\hline
启动加速度 $a_{1}$ & $4\Umsq$ \\
\hline
制动加速度 $a_{2}$ & $8\Umsq$ \\
\hline
直道最大速度 $v_{1}$ & $40\Ums$ \\
\hline
弯道最大速度 $v_{2}$ & $20\Ums$ \\
\hline
直道长度 $s$ & $218\Um$ \\
\hline
\end{tblr}
\end{table}

某同学是这样解的：要使摩托车所用时间最短，应先由静止加速到最大速度 $v_{1}=40\Ums$，然后再减速到 $v_{2}=20\Ums$，$t_{1}=\frac{v_{1}}{a_{1}}=\cdots\cdots$；$t_{2}=\frac{v_{1}-v_{2}}{a_{2}}=\cdots\cdots$；$t=t_{1}+t_{2}$。

你认为这位同学的解法是否合理？若合理，请完成计算；若不合理，请说明理由，并用你自己的方法算出正确结果。

%% answer:
\jdanswer{不正确。$11\Us$。}

%% memo:
\memoanswer{
该同学的解法不正确，因为摩托车必须在 $218\Um$ 的直道上完成变速运动过程，但按照该同学的解法，

$t_{1}=\frac{v_{1}}{a_{1}}=\frac{40}{4}\Us=10\Us$，$t_{2}=\frac{v_{1}-v_{2}}{a_{2}}=\frac{40-20}{8}\Us=2.5\Us$，

则 $t=t_{1}+t_{2}=12.5\Us$，

摩托车的位移为

$s=s_{1}+s_{2}=\frac{1}{2}v_{1}t_{1}+\frac{1}{2}(v_{1}+v_{2})t_{2}=0.5\times40\times10+0.5(40+20)\times2.5=275\Um$，

已大于直道长度 $218\Um$。

正确的解法如下：摩托车在 $t_{1}$ 时间内加速到 $v_{m}$，再在 $t_{2}$ 时间内减速到 $v_{2}$，总位移 $s$ 为 $218\Um$，

$t_{1}=\frac{v_{m}}{a_{1}}$ ①

$t_{2}=\frac{v_{m}-v_{2}}{a_{2}}$ ②

$\frac{1}{2}v_{m}t_{1}+\frac{1}{2}(v_{m}+v_{2})t_{2}=s$ ③

由①，②，③式联立解得 $v_{m}=36\Ums$ ④

最短时间

$t=t_{1}+t_{2}=\frac{v_{m}}{a_{1}}+\frac{v_{m}-v_{2}}{a_{2}}=\frac{36}{4}+\frac{36-20}{8}=11\Us$ ⑤
}

\item
%% number: 21
%% paperName: 2006年上海高考第 21 题 12 分
%% typeId: 6
%% type: 计算题
%% chapter: 运动和力
%% point: 加速减速模型
%% method: 无
%% score: 12
%% degree: 450
%% duplicateId: 0
%% body:
质量为 $10\Ukg$ 的物体在 $F=200\UN$ 的水平推力作用下，从粗糙斜面的底端由静止开始沿斜面运动，斜面固定不动，与水平地面的夹角 $\theta=37^{\circ}$。力 $F$ 作用 $2\Us$ 后撤去，物体在斜面上继续上滑了 $1.25\Us$ 后，速度减为零。求物体与斜面间的动摩擦因数 $\mu$ 和物体的总位移 $s$。（已知 $\sin37^{\circ}=0.6$，$\cos37^{\circ}=0.8$，$g=10\Umsq$）

\onepicture[w=5.0cm, align=right]{\includesvg{figs/21}}

%% answer:
\jdanswer{$\mu=0.25$，$s=16.25\Um$}

%% memo:
\memoanswer{
物体的整个运动分为两部分，设加速时的加速度为 $a_{1}$，减速时的加速度为 $a_{2}$，撤去力 $F$ 瞬间物体的速度为 $v$，则

$v=a_{1}t_{1}$，$0=v-a_{2}t_{2}$，

得 $a_{1}t_{1}=a_{2}t_{2}$ 或 $2a_{1}=1.25a_{2}$ ①

$a_{1}=\frac{F\cos\theta-mg\sin\theta-\mu(F\sin\theta+mg\cos\theta)}{m}$ ②

$a_{2}=\frac{mg\sin\theta+\mu mg\cos\theta}{m}$ ③

由①，②，③式联立解得：$\mu=0.25$ ④

代入②，③得：$a_{1}=5\Umsq$，$a_{2}=8\Umsq$。

物体的总位移为

$s=\frac{1}{2}a_{1}t_{1}^{2}+\frac{1}{2}a_{2}t_{2}^{2}$ ⑤

$=0.5\times5\times2^{2}+0.5\times8\times1.25^{2}=16.25\Um$ ⑥
}

\item
%% number: 22
%% paperName: 2006年上海高考第 22 题 14 分
%% typeId: 6
%% type: 计算题
%% chapter: 电磁感应
%% point: 线圈过有界磁场
%% method: 无
%% score: 14
%% degree: 450
%% duplicateId: 0
%% body:
如图所示，将边长为 $a$、质量为 $m$、电阻为 $R$ 的正方形导线框竖直向上抛出，穿过宽度为 $b$、磁感应强度为 $B$ 的匀强磁场，磁场的方向垂直纸面向里。线框向上离开磁场时的速度刚好是进入磁场时速度的一半，线框离开磁场后继续上升一段高度，然后落下并匀速进入磁场。整个运动过程中始终存在着大小恒定的空气阻力 $f$ 且线框不发生转动。求：

\begin{enumerate}
	\item
	线框在下落阶段匀速进入磁场时的速度 $v_{2}$；
	\item
	线框在上升阶段刚离开磁场时的速度 $v_{1}$；
	\item
	线框在上升阶段通过磁场过程中产生的焦耳热 $Q$。
\end{enumerate}

\onepicture[w=5.5cm, align=right]{\includesvg{figs/22}}

%% answer:
\jdanswer{
\begin{enumerate}
	\item
	$v_{2}=\frac{(mg-f)R}{B^{2}a^{2}}$
	\item
	$v_{1}=\frac{R\sqrt{(mg)^{2}-f^{2}}}{B^{2}a^{2}}$
	\item
	$Q=\frac{3}{2}m[(mg)^{2}-f^{2}]\frac{R^{2}}{B^{4}a^{4}}-(mg+f)(a+b)$
\end{enumerate}
}

%% memo:
\memoanswer{
（1）线框在下落阶段匀速进入磁场瞬间，由受力平衡得

$mg=f+\frac{B^{2}a^{2}v_{2}}{R}$ ①

解得 $v_{2}=\frac{(mg-f)R}{B^{2}a^{2}}$ ②

（2）线框从离开磁场至上升到最高点的过程，由动能定理得

$(mg+f)h=\frac{1}{2}mv_{1}^{2}$ ③

线框从最高点回落至磁场的过程

$(mg-f)h=\frac{1}{2}mv_{2}^{2}$ ④

③、④式联立解得

$v_{1}=\sqrt{\frac{mg+f}{mg-f}}v_{2}$ ⑤

$=\frac{R\sqrt{(mg)^{2}-f^{2}}}{B^{2}a^{2}}$ ⑥

（3）线框在向上通过磁场过程中，由动能定理得

$\frac{1}{2}mv_{0}^{2}-\frac{1}{2}mv_{1}^{2}=Q+(mg+f)(a+b)$ ⑦

又 $v_{0}=2v_{1}$，联立解得

$Q=\frac{3}{2}m[(mg)^{2}-f^{2}]\frac{R^{2}}{B^{4}a^{4}}-(mg+f)(a+b)$ ⑧
}

\item
%% number: 23
%% paperName: 2006年上海高考第 23 题 14 分
%% typeId: 6
%% type: 计算题
%% chapter: 电场
%% point: 电势能电势差电场力做功
%% method: 无
%% score: 14
%% degree: 250
%% duplicateId: 0
%% body:
电偶极子模型是指电量为 $q$、相距为 $l$ 的一对正负点电荷组成的电结构，O 是中点，电偶极子的方向为从负电荷指向正电荷，用图\ref{2006上海23a}所示的矢量表示。科学家在描述某类物质的电性质时，认为物质是由大量的电偶极子组成的，平时由于电偶极子的排列方向杂乱无章，因而该物质不显示带电的特性。当加上外电场后，电偶极子绕其中心转动，最后都趋向于沿外电场方向排列，从而使物质中的合电场发生变化。

\begin{enumerate}
	\item
	如图\ref{2006上海23b}所示，有一电偶极子放置在电场强度为 $E_{0}$ 的匀强外电场中，若电偶极子的方向与外电场方向的夹角为 $\theta$，求作用在电偶极子上的电场力绕 O 点的力矩；
	\item
	求图\ref{2006上海23b}中的电偶极子在力矩的作用下转动到外电场方向的过程中，电场力所做的功；
	\item
	求电偶极子在外电场中处于力矩平衡时，其方向与外电场方向夹角的可能值及相应的电势能；
	\item
	现考察物质中的三个电偶极子，其中心在一条直线上，初始时刻如图\ref{2006上海23c}排列，它们相互间隔距离恰等于 $l$。加上外电场 $E_{0}$ 后，三个电偶极子转到外电场方向，若在图中 A 点处引入一电量为 $+q_{0}$ 的点电荷（$q_{0}$ 很小，不影响周围电场的分布），求该点电荷所受电场力的大小。
\end{enumerate}

\threepicture[h=2.0cm, align=right, labelA=2006上海23a, labelB=2006上海23b, labelC=2006上海23c]{figs/23a.png}{figs/23b.png}{figs/23c.png}

%% answer:
\jdanswer{
\begin{enumerate}
	\item
	$M=qE_{0}l\sin\theta$
	\item
	$W=qE_{0}l(1-\cos\theta)$
	\item
	$\theta_{1}=0$ 时，$E_{\mathrm{p}1}=-qE_{0}l$；$\theta_{2}=\pi$ 时，$E_{\mathrm{p}2}=qE_{0}l$
	\item
	$F=q_{0}E_{0}-\frac{8kq_{0}q}{9l^{2}}$
\end{enumerate}
}

%% memo:
\memoanswer{
（1）由题意知，$+q$ 所受电场力矩为

$M_{1}=qE_{0}\frac{l}{2}\sin\theta$ ①

$-q$ 所受电场力矩为

$M_{2}=qE_{0}\frac{l}{2}\sin\theta$ ②

电偶极子所受力矩为

$M=M_{1}+M_{2}=qE_{0}l\sin\theta$ ③

（2）电场力对 $+q$ 做功

$W_{1}=qE_{0}\frac{l}{2}(1-\cos\theta)$ ④

电场力对 $-q$ 做功

$W_{2}=qE_{0}\frac{l}{2}(1-\cos\theta)$ ⑤

电场力对电偶极子做功

$W=W_{1}+W_{2}=qE_{0}l(1-\cos\theta)$ ⑥

（3）由③式 $M=qE_{0}l\sin\theta=0$ ⑦

得 $\theta_{1}=0$ 或 $\theta_{2}=\pi$。

$\theta_{1}=0$ 时，设此时点电荷所在位置的电势为 $\varphi$，电偶极子的电势能为

$E_{\mathrm{p}1}=-q\varphi+q(\varphi-E_{0}l)=-qE_{0}l$，

$\theta_{2}=\pi$ 时，电偶极子的电势能为

$E_{\mathrm{p}2}=-q\varphi+q(\varphi+E_{0}l)=qE_{0}l$ ⑧

（4）三个电偶极子沿电场方向排列，中间的正负电荷相互抵消，两端的正电荷作用在 A 点处点电荷 $q_{0}$ 上的电场大小均为 $\frac{kq_{0}q}{\left(\frac{3}{2}l\right)^{2}}$，方向均与外电场相反，三个电偶极子作用在 $q_{0}$ 上的电场力为

$F'=\frac{2kq_{0}q}{\left(\frac{3}{2}l\right)^{2}}=\frac{8kq_{0}q}{9l^{2}}$ ⑨

外电场作用在 $q_{0}$ 上的电场力为

$F=q_{0}E_{0}$，

$q_{0}$ 所受的电场力为

$F=F_{0}-F'=q_{0}E_{0}-\frac{8kq_{0}q}{9l^{2}}$ ⑩
}

\end{enumerate}

\end{document}
